\section{Design}
\label{sec:design}

\subsection{Overview}
\label{sec:design-overview}

\Cref{fig:architecture} shows the control loop. Context observations enter a
gateway-side belief filter that maintains a distribution over five latent threat
states and a vector of modality-trust weights. The belief and the current
operation class are passed to an admissible-set constructor, which applies the
unconditional floors I1/I2 and the conditional restriction I3. A constrained
scorer ranks the surviving actions by expected utility, subject to hysteresis
and a switching margin, and the policy enforcement point executes the selected
action. A precomputed fail-safe action is available on the bypass path for I4.
The key architectural property is that the belief filter can influence
\emph{which} safe action is chosen, but not \emph{whether} the action is safe.

\begin{figure}[t]
\centering
\resizebox{\textwidth}{!}{%
\begin{tikzpicture}[
  font=\scriptsize,
  node distance=5mm and 7mm,
  box/.style={draw, rounded corners=2pt, align=center, minimum height=9mm,
              text width=23mm, inner sep=2pt},
  smallbox/.style={draw, rounded corners=2pt, align=center, minimum height=7mm,
                   text width=19mm, inner sep=2pt, fill=black!4},
  arr/.style={-{Latex[length=1.6mm]}, thick},
  darrow/.style={-{Latex[length=1.6mm]}, thick, dashed},
]
\node[box] (ctx) {Device, network, security and cross-node context};
\node[box, right=of ctx] (belief) {Belief filter\\HMM $+$ trust weights};
\node[box, right=of belief] (safe) {Admissible set\\I1, I2, I3};
\node[box, right=of safe] (score) {Constrained scorer\\utility, dwell, margin};
\node[box, right=of score] (pep) {PEP\\action or deny};
\node[box, below=8mm of pep] (proc) {Physical process};
\node[smallbox, below=8mm of safe] (fail) {Fail-safe action\\I4 (per operation)};

\draw[arr] (ctx) -- (belief);
\draw[arr] (belief) -- (safe);
\draw[arr] (safe) -- (score);
\draw[arr] (score) -- (pep);
\draw[arr] (pep) -- (proc);
\draw[darrow] (proc.west) -| ++(-8mm,-4mm) -| (ctx.south)
  node[pos=0.25, below, font=\tiny] {feedback};
\draw[arr] (fail) -- (pep);
\draw[darrow] (belief.south) |- (fail.west);
\node[draw, dashed, rounded corners=3pt, fit=(belief)(safe)(score), inner sep=2mm,
      label={[font=\tiny]above:edge gateway (trusted PDP)}] {};
\node[draw, dashed, rounded corners=3pt, fit=(pep)(proc), inner sep=2mm,
      label={[font=\tiny]below:node (PEP)}] {};
\end{tikzpicture}}
\caption{Control loop. The belief filter and scorer run on the gateway; the PEP
runs on the node. I1/I2 are enforced in the admissible-set constructor and do
not depend on the belief. The fail-safe path is precomputed per operation class.}
\label{fig:architecture}
\end{figure}

\subsection{Belief over latent threat states}
\label{sec:belief}

Let $\Theta=\{\text{benign},\text{recon},\text{DoS},\text{Byzantine},\text{compromised}\}$
and let $b_t(\btheta)=\Prob(\theta_t=\btheta\mid\bx_{1:t})$ be the defender's
belief. The latent state evolves as a stationary Markov chain with transition
matrix $T$, where $\operatorname{diag}(T)=0.82$ and the remaining mass is spread
uniformly over the other four states. The persistence is deliberate: threat
states are episodes, not independent draws, and a filter that ignores persistence
over-reacts to single observations.

For each modality $m$ with coordinate set $I_m$, the emission model is a product
of univariate Gaussians,
\begin{equation}
  L_m(\bx_t\mid\btheta)
  =\prod_{j\in I_m}\frac{1}{\sigma_{\btheta,j}}
    \exp\!\left[-\frac{(x_{t,j}-\mu_{\btheta,j})^2}{2\sigma_{\btheta,j}^2}\right],
  \label{eq:emission}
\end{equation}
with state- and feature-specific means and standard deviations. In this
simulator the likelihood is evaluated at the same parameters that generate the
observations, so the defender is correctly specified by construction and the
calibration and detection results are an oracle-model upper bound, not an
estimate for a misspecified defender. The classes still overlap by design: a DoS episode and a heavily loaded benign node can produce
similar device telemetry, and a Byzantine node can look almost benign except for
cross-node dispersion. The model is intentionally not tuned so that a single
feature separates the classes.

\paragraph{Trust weights.} For each modality we form a \emph{modality posterior}
$q_m(\btheta)\propto \pi_0(\btheta)L_m(\bx_t\mid\btheta)$ using a fixed reference
prior $\pi_0$, and compare it with the device posterior $q_{\mathrm{dev}}$ by
total variation:
\begin{equation}
  \tau_m = \tfrac12\lVert q_m-q_{\mathrm{dev}}\rVert_1 .
  \label{eq:tv}
\end{equation}
The weight is a logistic gate
\begin{equation}
  \tilde w_m = \frac{1}{1+\exp[(\tau_m-\tau_{\mathrm{TV}})/s_{\mathrm{TV}}]},
  \qquad \tilde w_{\mathrm{dev}}=\tilde w_{\mathrm{cross}}=1 ,
  \label{eq:weight}
\end{equation}
so a modality whose posterior has drifted far from the attestation anchor is
discarded, while the device anchor and the dispersion signal are always kept.
The weight is updated asymmetrically,
\begin{equation}
  w_{m,t}=
  \begin{cases}
    0.2\,w_{m,t-1}+0.8\,\tilde w_m, & \tilde w_m<w_{m,t-1},\\
    0.9\,w_{m,t-1}+0.1\,\tilde w_m, & \text{otherwise},
  \end{cases}
  \label{eq:asym}
\end{equation}
which distrusts quickly (80\,\% of the new evidence) and restores slowly
(10\,\%). Slow restoration prevents a brief clean observation from undoing a
sustained inconsistency.

\paragraph{Fusion.} The fused belief is a forward HMM update with weighted
modality likelihoods:
\begin{equation}
  b_t(\btheta)=\frac{1}{Z_t}
  \Big[\sum_{\btheta'}T_{\btheta'\btheta}\,b_{t-1}(\btheta')\Big]
  \prod_{m\in\mathcal{M}}L_m(\bx_t\mid\btheta)^{w_{m,t}},
  \label{eq:fusion}
\end{equation}
where $Z_t$ normalises the distribution. A separate \emph{uniform-fusion} belief
uses $w_{m,t}=1$ for all modalities; it is the ablation that isolates the value
of trust weighting. The attack probability is
$\attack_t=b_t(\text{DoS})+b_t(\text{Byzantine})+b_t(\text{compromised})$;
reconnaissance is tracked but excluded from the serious-attack mass because it
does not by itself justify a hard restriction.

\paragraph{Consistency score.} For monitoring and for I3 we also need a
belief-free signal. Let $d_t$ be the cross-node dispersion and let
\begin{equation}
  s_t=\max\Big\{\max_{m}\tau_m,\;
  \operatorname{clip}\!\left(\frac{d_t-d_0}{d_1},0,1\right)\Big\},
  \label{eq:consistency}
\end{equation}
with $d_0,d_1$ calibrated from benign traffic. The first term catches a modality
that disagrees with the device anchor; the second catches a colluding minority
that shifts the cross-node spread even when every individual modality looks
plausible. This score is used for detection metrics and as a diagnostic; it is
not the safety mechanism.

\subsection{Action lattice, detection and utility}
\label{sec:utility}

An action $a=(a_{\mathrm{auth}},a_{\mathrm{cry}},a_{\mathrm{ids}},a_{\mathrm{rate}},a_{\mathrm{ttl}})$
is a point in the 243-element lattice. For a malicious request, the probability
that the action detects it is the complement of the probability that every layer
misses:
\begin{equation}
  p_{\mathrm{det}}(\btheta,a)=1-\prod_{\ell}\big(1-p_{\ell}(\btheta,a_\ell)\big),
  \label{eq:detect}
\end{equation}
with base validation, authentication, inspection and rate-limit terms. The
state-dependent multipliers make inspection less effective against a Byzantine
or compromised node (the adversary knows the signatures), and the authentication
term rewards attestation over bearer tokens. For a legitimate request, the
false-positive probability $p_{\mathrm{fp}}(a)$ is the union of the same layers'
false-alarm probabilities.

Let $V_{\op}$ be the value of a granted legitimate request, $D_{\op}$ the damage
of a granted malicious one, and $C_{\op}$ the cost of denying a legitimate
request. With $p_{\mathrm{mal}}(\btheta)=\Prob(\text{request is malicious}\mid\btheta)$,
the expected utility of action $a$ for operation $\op$ in state $\btheta$ is
\begin{equation}
\begin{aligned}
  U_{\op}(\btheta,a)={}&
  p_{\mathrm{mal}}(\btheta)\big[-(1-p_{\mathrm{det}}(\btheta,a))D_{\op}\big]\\
  &+(1-p_{\mathrm{mal}}(\btheta))\big[V_{\op}(1-p_{\mathrm{fp}}(a))-C_{\op}p_{\mathrm{fp}}(a)\big].
\end{aligned}
  \label{eq:utility}
\end{equation}
Safety-critical operations are never malicious in the model, so
$p_{\mathrm{mal}}=0$ for them; the design question there is availability, not
malice. The action score adds the systems cost:
\begin{equation}
  S_{\op}(a;b_t)=b_t^{\top}U_{\op}(\cdot,a)
  -\lambda_L\ell(a)-\lambda_E e(a),
  \label{eq:score}
\end{equation}
where $\ell(a)$ and $e(a)$ are latency and energy, and $\lambda_L,\lambda_E$
convert them to utility units. The realised utility used for reporting adds one
further term, $\rho\,\ind{a\not\models\mathrm{I1}\land a\neq\textsc{Deny}}$
with $\rho=6.0$, which prices a granted action that falls below the floor. This
term is not part of the controller's score---the controller cannot select such
an action---and it exists so that the utility metric does not reward a violation
that the invariants forbid. Denying has score
$-(1-p_{\mathrm{mal}}^{\mathrm{req}})C_{\op}$ with
$p_{\mathrm{mal}}^{\mathrm{req}}=(1-b_t(\text{benign}))P_{\mathrm{malicious}}$:
denying is cheap when the request is likely malicious, but not free when
legitimate traffic is still arriving. This request-level term is important; an
epoch-level denial heuristic over-denies mixed traffic.

\subsection{Admissible set and constrained objective}
\label{sec:admissible}

For operation $\op$ and belief $b$, the admissible set is
\begin{equation}
  \Asafe(\op,b)=\{a\in\Aspace:\; a\models \mathrm{I1}(\op)\land \mathrm{I2}(\op)
  \land \mathrm{I3}(\op,b)\},
  \label{eq:admissible}
\end{equation}
where I1/I2 are the unconditional floors and I3 is the belief-conditional
restriction. Deny is a separate action, admissible only when
$\op\neq\text{safety}$. The controller solves, at each request,
\begin{equation}
  a_t^{\star}=\argmax_{a\in\Asafe(\op,\attack_t)\cup\mathcal{D}(\op)}
  \Big[S_{\op}(a;b_t)-\gamma\,\ind{a\neq a_{t-1}}\Big],
  \label{eq:objective}
\end{equation}
where $\mathcal{D}(\op)$ is the deny singleton for non-safety operations and
$\gamma$ is the switching penalty. The objective is myopic by design: the belief
already carries the temporal state, and a receding-horizon planner would add
latency without changing the safety argument.

\begin{proposition}[Invariant dominance and non-emptiness]
\label{prop:dominance}
For every operation $\op$ and every belief $b\in\Delta(\Theta)$, the set
$\Asafe(\op,b)$ is non-empty. It contains the maximal posture
$a_{\max}=(\text{attestation},\text{Ascon},\text{deep},\text{10/s},\text{10\,s})$,
which satisfies I1, I2 and I3. Consequently, any action returned by
\eqref{eq:objective} is in $\Asafe$, and the fail-safe fallback in I4 is always
defined.
\end{proposition}

\begin{proof}
At $a_{\max}$ the authentication level is 3, the AEAD level is 2, the inspection
level is 2, the rate-limit level is 2 and the freshness level is 2. I1 requires
authentication $\ge 2$, AEAD, freshness $\ge 1$ and inspection $\ge 1$; I2
requires authentication $\ge 2$, inspection $\ge 2$ and freshness $\ge 1$; I3
can only add lower bounds at or below these levels. Hence $a_{\max}$ satisfies
all three for every $\op$ and every $b$. Selection is restricted to $\Asafe$ by
\eqref{eq:objective}; if the set is empty or the controller times out, the PEP
executes $a_{\max}$ (or the operation-specific safe action), which is also
admissible.
\end{proof}

\begin{proposition}[Structural safety for any belief]
\label{prop:structural}
For any belief sequence $\{b_t\}$, any request that is granted by the controller
uses an action that satisfies I1, and, if the request is safety-critical, I2.
Equivalently, the invariant-violation rate is zero by construction, not by
estimation accuracy.
\end{proposition}

\begin{proof}
A granted request is either (i) the output of \eqref{eq:objective}, which lies
in the admissible set $\Asafe$ for that operation and belief and therefore
satisfies I1 and, for safety-critical operations, I2, or (ii) the I4 fallback,
which is $a_{\max}$ and satisfies the same floors. A denied request is not
granted and therefore cannot violate the floors. The claim holds for an
adversarially chosen belief because no step uses a correctness assumption about
$b_t$.
\end{proof}

\begin{proposition}[Safety liveness]
\label{prop:liveness}
For $\op=\text{safety}$, deny is not in the admissible action set for any belief,
and $\Asafe(\text{safety},b)$ is non-empty. A safety-critical request is
therefore never denied by policy; it can only be denied by a stochastic
mechanism failure (authentication or inspection false positive).
\end{proposition}

\begin{proof}
I4 excludes deny for safety, and \cref{prop:dominance} gives a non-empty
admissible set for every belief. The controller must therefore select an action
from $\Asafe(\text{safety},b)$, and the PEP grants it subject only to the
mechanism's false-positive probability.
\end{proof}

\subsection{Fail-safe enforcement algorithm}
\label{sec:algorithm}

\Cref{alg:controller} combines the belief update, admissible-set construction,
scoring, hysteresis and fallback. The hysteresis state is per operation class:
the controller remembers the previous security level and refuses to de-escalate
while the risk is inside the band $[\tau_{\mathrm{down}},\tau_{\mathrm{up}})$
and the low-risk streak is shorter than the dwell time. A lateral or
de-escalating move must beat the incumbent by a margin $m$; an escalation is
immediate. This asymmetry is intentional: safety is allowed to react quickly,
availability is not.

\begin{algorithm}[t]
\caption{Fail-safe context-adaptive enforcement (per epoch)}
\label{alg:controller}
\begin{algorithmic}[1]
\Require observation $\bx_t$; requests $\mathcal{R}_t$; previous state
$\{(a_{\op},s_{\op})\}_{\op}$; low-risk streak $\ell$
\Ensure actions for all requests
\State $q_m\propto\pi_0\odot L_m(\bx_t)$ for each modality $m$
\State $\tau_m\gets\tfrac12\lVert q_m-q_{\mathrm{dev}}\rVert_1$;\quad
       $\tilde w_m\gets\sigma((\tau_{\mathrm{TV}}-\tau_m)/s_{\mathrm{TV}})$
\State $w_t\gets\textsc{AsymUpdate}(w_{t-1},\tilde w)$ \Comment{Eq.~\eqref{eq:asym}}
\State $b_t\gets\textsc{Normalise}\big[(T^{\top}b_{t-1})\odot
       \exp(\sum_m w_{m,t}\log L_m)\big]$ \Comment{Eq.~\eqref{eq:fusion}}
\State $\attack_t\gets b_t(\text{DoS})+b_t(\text{Byz})+b_t(\text{comp})$
\State $\ell\gets\ell+1$ if $\attack_t<\tau_{\mathrm{down}}$ else $0$
\ForAll{$(\op,\ldots)\in\mathcal{R}_t$}
  \State $\Asafe\gets\{a:a\models\mathrm{I1}(\op)\land\mathrm{I2}(\op)\land
         \mathrm{I3}(\op,\attack_t)\}$
  \State $\mathcal{C}\gets\Asafe$;\quad
         if $\op\neq\text{safety}$ then $\mathcal{C}\gets\mathcal{C}\cup\{\textsc{Deny}\}$
  \If{$\mathcal{C}=\emptyset$}
     \State $a^{\star}\gets\textsc{SafeFallback}(\op)$ \Comment{I4}
  \Else
     \State $a^{\star}\gets\argmax_{a\in\mathcal{C}}S_{\op}(a;b_t)-\gamma\ind{a\neq a_{\op}}$
     \State apply dwell floor and margin $m$ to $a^{\star}$
  \EndIf
  \State enforce $a^{\star}$ (or deny if $a^{\star}=\textsc{Deny}$)
  \State update $(a_{\op},s_{\op})$
\EndFor
\end{algorithmic}
\end{algorithm}

\paragraph{Complexity.} The belief update is $O(|\Theta|^2+|\mathcal{M}||\Theta|)$
per epoch, the admissible-set construction is $O(|\Aspace|)$, and scoring is
$O(|\Aspace|)$ per request. With $|\Aspace|=243$ and $|\Theta|=5$, the whole
decision is a few thousand floating-point operations; \cref{sec:evaluation}
shows that it is dominated by the cryptographic mechanisms it selects, not by
the controller.

\paragraph{Why these choices.}
The device modality is the anchor because it is the only context source that is
attested by the node itself; an attacker who controls the network can make the
network and security modalities agree with each other, but not with a value
signed by the device. The cross-node dispersion is kept at full weight because a
colluding minority cannot hide the spread it creates without controlling the
majority. This anchor is an explicit trust assumption: the \texttt{tee}
scenario removes it, and the belief is expected to fail there. Asymmetric trust is a deliberate bias: a false distrust costs some
utility, while a false trust can cost a physical action. The utility margin is
the analogue of a Schmitt trigger; it converts small belief fluctuations into
stability without hiding large ones.
